\documentclass[10pt,a4paper]{amsart}
\usepackage[margin=19mm]{geometry}
\usepackage{amsmath,amssymb,mathtools}
\usepackage[scr=rsfs,cal=euler]{mathalfa}
\usepackage{xfrac}
\usepackage[unicode,hidelinks,bookmarks=false]{hyperref}
\newcommand{\ie}{i.e.\ }
\newcommand{\cf}{cf.\ }
\newcommand{\resp}{respectively\ }
\newcommand{\Subsection}{\subsection}
\providecommand{\nobreakdash}{\nobreak\hbox{-}}
\setlength{\emergencystretch}{2em}
\multlinegap0pt
\usepackage{amssymb}
\usepackage{mathtools}
\usepackage{bm}
\usepackage[shortlabels]{enumitem}
\usepackage{xcolor}
\newtheorem{lemma}{Lemma}
\providecommand{\C}{\mathbb{C}}
\providecommand{\Hphys}{\mathcal H_{\mathrm{phys}}}
\providecommand{\Hzcurl}{H_0(\operatorname{curl})}
\providecommand{\OmegaLa}{\Omega_{L,a}}
\providecommand{\R}{\mathbb{R}}
\providecommand{\VPEC}{\mathcal V_{\mathrm{PEC}}}
\providecommand{\Vphys}{\mathcal V_{\mathrm{phys}}}
\providecommand{\Z}{\mathbb{Z}}
\providecommand{\abs}[1]{\left\lvert #1\right\rvert}
\providecommand{\hzero}{h_0}
\providecommand{\ii}{\mathrm{i}}
\providecommand{\norm}[1]{\left\lVert #1\right\rVert}
\providecommand{\paren}[1]{\left(#1\right)}
\providecommand{\set}[1]{\left\{#1\right\}}
\title{A Maxwell Quadratic-Form Representation of the Parallel-Plate Casimir Trace from Codimension-Three Riesz Reduction}
\author{Irshadullah Khan, Bilal Khan}
\date{Source: arXiv:2605.20230v1 (CC BY 4.0)}
\begin{document}
\maketitle

\noindent\emph{Source and licence.} A Maxwell Quadratic-Form Representation of the Parallel-Plate Casimir Trace from Codimension-Three Riesz Reduction. Version arXiv:2605.20230v1 (CC BY 4.0), distributed by arXiv under the Creative Commons Attribution 4.0 International licence, http://creativecommons.org/licenses/by/4.0/, as recorded in the arXiv licence metadata for this exact version and shown on the abstract page https://arxiv.org/abs/2605.20230v1. Full licence notice: \texttt{licenses/arxiv-2605.20230-CC-BY-4.0.txt}.\par\smallskip

\noindent\textit{Editorial scope.} Counted unit: the article's lemma lem:app-fourier-characterization (source0.tex lines 3143--3167). The article's complete original proof of it is delivered with it (source0.tex lines 3169--3304, one \texttt{proof} environment of 136 lines). No mathematics is written, repaired or re-derived: the statement and the proof are the article's own lines, in the article's own notation, and no retained line is substituted or re-flowed.  Retained uncounted as the source-local context the proof needs: lines 3084--3095; lines 3097--3104; lines 3121--3129; lines 3131--3134.  Nothing was omitted from any retained span, so the package carries no omission record.  No retained line contains a citation, so the wrapper carries no bibliography and no bibliography entry.  No retained line contains a figure environment, an image-inclusion command or drawing code, so no image is delivered and the package declares no asset dependency.  Editorial formatting only, in addition: an AMS \texttt{amsart} preamble that restates the article's own declarations and macros used by the retained lines, namely the environments \texttt{lemma} on separate counters, together with the standard packages those lines need.  The retained environments print in this document as Lemma 1 in order of appearance, whereas the article prints the same items with the numbering of its own full document.  The complete native source of the article as distributed in the arXiv e-print for this exact version is bundled unchanged as \texttt{sources/200979/source0.tex}.
\begin{align}
    E_{\parallel}(r,z)
    &=
    \sum_{m\in\Z^2}\sum_{n\ge1}
    a_{m,n}\,\phi_m(r)s_n(z),
    \qquad a_{m,n}\in\C^2,                                  \label{eq:app-mixed-tangential}\\
    E_z(r,z)
    &=
    \sum_{m\in\Z^2}\sum_{n\ge0}
    b_{m,n}\,\phi_m(r)c_n(z),
    \qquad b_{m,n}\in\C .                                     \label{eq:app-mixed-normal}
\end{align}

\begin{equation}
\label{eq:app-parseval-L2}
    \norm{E}_{L^2(\OmegaLa)}^2
    =
    \sum_{m\in\Z^2}\sum_{n\ge1}\abs{a_{m,n}}^2
    +
    \sum_{m\in\Z^2}\sum_{n\ge0}\abs{b_{m,n}}^2 .
\end{equation}

\begin{equation}
\label{eq:app-CPEC-summability}
    \sum_{m\in\Z^2}\sum_{n\ge1}
    (1+\lambda_{m,n})
    \paren{\abs{a_{m,n}}^2+\abs{b_{m,n}}^2}
    +
    \sum_{m\in\Z^2}(1+\abs{k_m}^2)\abs{b_{m,0}}^2
    <\infty
\end{equation}

\begin{equation}
\label{eq:app-div-constraint}
    \ii k_m\cdot a_{m,n}-\nu_n b_{m,n}=0 .
\end{equation}

\begin{lemma}[Detailed Fourier characterization]
\label{lem:app-fourier-characterization}
The reconstruction map \eqref{eq:app-mixed-tangential}--\eqref{eq:app-mixed-normal}
identifies \(\mathfrak C_{\mathrm{PEC}}\) with
\[
    \VPEC=
    \set{E\in\Hzcurl(\OmegaLa;\C^3):\nabla\cdot E=0
    \text{ in }\mathcal D'(\OmegaLa)} .
\]
For \(E\in\VPEC\), the curl norm is
\begin{equation}
\label{eq:app-curl-norm-diagonal}
    \norm{\nabla\times E}_{L^2}^2
    =
    \sum_{m\in\Z^2}\sum_{n\ge1}
    \lambda_{m,n}
    \paren{\abs{a_{m,n}}^2+\abs{b_{m,n}}^2}
    +
    \sum_{m\in\Z^2}\abs{k_m}^2\abs{b_{m,0}}^2 .
\end{equation}
Moreover, imposing \(b_{0,0}=0\) identifies \(\mathfrak C_{\mathrm{phys}}\) with
\[
    \Vphys=\VPEC\cap\Hphys .
\]
\end{lemma}

\begin{proof}
We split the proof into four elementary steps.

\emph{Step 1: divergence coefficients.}
For a finite trigonometric field of the form
\eqref{eq:app-mixed-tangential}--\eqref{eq:app-mixed-normal}, direct
differentiation gives
\begin{equation}
\label{eq:app-div-expansion}
    \nabla\cdot E
    =
    \sum_{m\in\Z^2}\sum_{n\ge1}
    \paren{\ii k_m\cdot a_{m,n}-\nu_n b_{m,n}}
    \phi_m(r)s_n(z).
\end{equation}
There is no contribution from \(b_{m,0}\), since \(c_0'(z)=0\).  Thus the
constraint \eqref{eq:app-div-constraint} is exactly the Fourier form of
\(\nabla\cdot E=0\).

\emph{Step 2: curl coefficients and the finite-dimensional identity.}
For \(n\ge1\), the curl coefficients are
\begin{align}
\label{eq:app-curl-coeffs-positive-n}
    C^x_{m,n}&=\ii k_{m,y}b_{m,n}-\nu_n a^y_{m,n},\nonumber\\
    C^y_{m,n}&=\nu_n a^x_{m,n}-\ii k_{m,x}b_{m,n},\\
    C^z_{m,n}&=\ii\paren{k_{m,x}a^y_{m,n}-k_{m,y}a^x_{m,n}}.\nonumber
\end{align}
For the normal branch \(n=0\),
\begin{equation}
\label{eq:app-curl-coeffs-zero-n}
    C^x_{m,0}=\ii k_{m,y}b_{m,0},
    \qquad
    C^y_{m,0}=-\ii k_{m,x}b_{m,0},
    \qquad
    C^z_{m,0}=0 .
\end{equation}
For \(n\ge1\), let
\[
    D_{m,n}:=\ii k_m\cdot a_{m,n}-\nu_n b_{m,n} .
\]
A direct expansion gives the identity
\begin{align}
\label{eq:app-curl-div-identity}
    &\abs{C^x_{m,n}}^2+
     \abs{C^y_{m,n}}^2+
     \abs{C^z_{m,n}}^2+
     \abs{D_{m,n}}^2        \\
    &\hspace{3cm}=
     \lambda_{m,n}
     \paren{\abs{a_{m,n}}^2+\abs{b_{m,n}}^2} .\nonumber
\end{align}
One way to check the cancellation is to observe that the terms involving
\(b_{m,n}\overline{a^x_{m,n}}\) and \(b_{m,n}\overline{a^y_{m,n}}\) occur with
opposite signs in the curl part and in \(\abs{D_{m,n}}^2\).  The remaining
terms use the two-dimensional identity
\[
    \abs{k_m\cdot a_{m,n}}^2+
    \abs{k_{m,x}a^y_{m,n}-k_{m,y}a^x_{m,n}}^2
    =
    \abs{k_m}^2\abs{a_{m,n}}^2,
\]
valid for complex \(a_{m,n}\in\C^2\) and real \(k_m\in\R^2\).

\emph{Step 3: from coefficients to the form domain.}
Let \((a,b)\in\mathfrak C_{\mathrm{PEC}}\).  Let \(E^{(N)}\) be any increasing
sequence of finite rectangular truncations of
\eqref{eq:app-mixed-tangential}--\eqref{eq:app-mixed-normal}.  Each truncation is
smooth, lateral-periodic, and has tangential components built from sine modes,
so
\[
    n\times E^{(N)}=0
    \qquad\text{on }z=0,a .
\]
The coefficient constraint \eqref{eq:app-div-constraint} is preserved under
truncation, hence \(\nabla\cdot E^{(N)}=0\).  By
\eqref{eq:app-curl-div-identity} with \(D_{m,n}=0\) and by
\eqref{eq:app-curl-coeffs-zero-n}, the summability condition
\eqref{eq:app-CPEC-summability} implies that \(E^{(N)}\) is Cauchy in the
\(H(\operatorname{curl})\) graph norm.  Since \(\Hzcurl\) is the graph-norm
closure of smooth PEC fields, the limit belongs to \(\Hzcurl\).  The divergence
constraint passes to the limit distributionally.  Therefore the reconstructed
field belongs to \(\VPEC\).  This proves
\[
    \mathfrak C_{\mathrm{PEC}}\subset\VPEC .
\]

\emph{Step 4: from the form domain to coefficients.}
Conversely, let \(E\in\VPEC\).  Its mixed Fourier coefficients are defined by
\eqref{eq:app-mixed-tangential}--\eqref{eq:app-mixed-normal}.  Since
\(\nabla\cdot E=0\) in distributions, testing against
\(\overline{\phi_m}s_n\), \(n\ge1\), gives precisely
\eqref{eq:app-div-constraint}.  No endpoint term arises in this test because
\(s_n(0)=s_n(a)=0\).

The weak curl \(\nabla\times E\) is in \(L^2\).  Its Fourier coefficients are the
distributional derivatives displayed in
\eqref{eq:app-curl-coeffs-positive-n}--\eqref{eq:app-curl-coeffs-zero-n}.  These
formulas are obtained by differentiating finite partial sums, testing against
the mixed sine/cosine modes, and passing to the distributional limit.  Since
\(\nabla\times E\in L^2\), uniqueness of Fourier coefficients identifies the
\(L^2\)-Fourier coefficients of the weak curl with the displayed expressions.
Parseval
therefore gives
\[
    \sum\abs{C}^2<\infty .
\]
Using \eqref{eq:app-curl-div-identity} and the fact that
\(D_{m,n}=0\), we obtain
\[
    \sum_{m\in\Z^2}\sum_{n\ge1}
    \lambda_{m,n}
    \paren{\abs{a_{m,n}}^2+\abs{b_{m,n}}^2}
    +
    \sum_{m\in\Z^2}\abs{k_m}^2\abs{b_{m,0}}^2
    <\infty .
\]
Together with the \(L^2\) Parseval identity \eqref{eq:app-parseval-L2}, this is
exactly \eqref{eq:app-CPEC-summability}.  Hence the coefficient family of
\(E\) lies in \(\mathfrak C_{\mathrm{PEC}}\), and the curl norm is
\eqref{eq:app-curl-norm-diagonal}.  Thus
\[
    \VPEC\subset\mathfrak C_{\mathrm{PEC}} .
\]
The two inclusions prove the identification.

Finally, the mode \(b_{0,0}\) reconstructs the vector field
\[
    b_{0,0}\phi_0(r)c_0(z)e_z
    =
    b_{0,0}(L^2a)^{-1/2}e_z
    =
    b_{0,0}\hzero .
\]
Thus the condition \(b_{0,0}=0\) is exactly orthogonality to the static normal
mode \(\hzero\).  This proves the reduced statement.
\end{proof}

\end{document}
